\documentclass{article}

\usepackage{arxiv}

\usepackage[utf8]{inputenc} % allow utf-8 input
\usepackage[T1]{fontenc}    % use 8-bit T1 fonts
\usepackage{setspace}
\setstretch{1.2}          % improved readability (line spacing)

\usepackage{hyperref}       % hyperlinks
\usepackage{url}            % simple URL typesetting
\usepackage{booktabs}       % professional-quality tables
\usepackage{amsfonts}       % blackboard math symbols
\usepackage{amsmath,amssymb}
\usepackage{nicefrac}       % compact symbols for 1/2, etc.
\usepackage{microtype}      % microtypography
\usepackage{graphicx}
\graphicspath{ {./images/} }
\providecommand{\tightlist}{%
  \setlength{\itemsep}{0pt}\setlength{\parskip}{0pt}}

\title{Neocortex\\\large artificial consciousness with a surprise-weighted context system}

\author{
 Steven Enamakel \\
  Tiny Humans \\
  \texttt{enamakel@tinyhumans.ai} \\
}

\begin{document}
\maketitle
\begin{abstract}
{\small
\begin{minipage}{0.82\textwidth}
\noindent
Current Large Language Models (LLMs) do not have consciousness or intrinsic intellect; they generate text based on training data by modeling language patterns rather than experiencing understanding. Furthermore, current LLMs are limited by bounded context windows; increasing them requires exponentially more compute and memory. While RAGs (Retrieval Augmented Generation) improve factual access to external corpora, they are not smart enough in simulating an artificial consciousness.
\\\\
\noindent
The human brain is an incredibly power-efficient, parallel noise-cutting engine. Inspired by this, We design Neocortex, a contextual memory system that implements memory at inference time. It unifies coherence memory, a state event ledger, interaction memory and forgetting/surprise dynamics, to create a powerful, compute-efficient, artificial consciousness.
\end{minipage}
}

\end{abstract}

\section{Introduction}

Modern LLMs can reason impressively within a single prompt, but
building long-horizon systems remains hard. Despite steady progress in
reasoning, LLMs are still constrained by finite context windows and weak
persistence across interactions. Reasoning-oriented LLMs can allocate
more computation to multi-step deliberation, yet they remain highly
sensitive to what is placed into context, and naive RAG retrieval often
treats memory as a flat collection of text chunks. Simply scaling context
windows is expensive: attention and KV-cache costs grow rapidly with
sequence length, and longer prompts can reduce accuracy by injecting
irrelevant material, diluting key evidence, and increasing the chance
that the model attends to the wrong details. Practical long-horizon
systems therefore need memory mechanisms that select, structure, and
refresh context rather than simply making the window bigger.

In parallel, a growing ecosystem of memory tooling has emerged that is
already highly efficient at storing and retrieving information for LLM
applications. Systems such as MemGPT, mem0, and Supermemory maintain
external memory stores and retrieval policies that can prioritize
memories by time/recency, entities and user profiles, interaction
history, or task-specific reasoning needs. These approaches demonstrate
that the core bottleneck is often not storage, but selecting the right
facts at the right time under tight latency and token budgets.

The human brain faces an extreme noise problem: it continuously receives
high-volume, ambiguous, and often redundant signals, yet it filters and
compresses them into a small set of salient memories that influence
future behavior. Many current LLM memory systems can store and retrieve
efficiently, but they struggle to consistently infer and surface
\emph{importance}---what should be reinforced, what should be forgotten,
and what should be promoted to durable state---especially when signals
are subtle, long-range, or conflicting. Human memory offers a useful
conceptual alternative: it is reinforced through use, decays over time,
and is selectively updated when new experiences are surprising or
contradictory, motivating forgetting-aware updates, episodic
segmentation, graph-based retrieval, and persistent memory modules.

This paper introduces \textbf{Neocortex}, a product-level memory
architecture that unifies these strands into a single operational
pipeline. Neocortex continuously ingests new experience into structured
memory and recalls context through a blend of semantic relevance,
recency, interaction history, and surprise-weighted salience. At write
time, the system parses documents into chunks, extracts entities and
relations, persists them in graph/vector memory, and appends explicit
state transitions to an event ledger. At recall time, the system routes
queries either to broad semantic retrieval or to a deterministic state
resolver when the question targets ordered state transitions. Memory is
further modulated by reinforcement through access patterns and by
Ebbinghaus-style forgetting dynamics.

Our main claim is that long-term AI memory should be treated as a
\textbf{multi-store adaptive system} rather than a passive retrieval
index. In particular, efficient token use requires not only better
retrieval, but also mechanisms for forgetting noise, amplifying
informative novelty, and separating semantic context from ordered state.

\section{Related Work}

\subsection{Long-term memory for language
models}

Several recent systems address persistent memory in LLMs. MemoryBank
augments conversational agents with a long-term memory store and a
forgetting-aware update mechanism inspired by the Ebbinghaus forgetting
curve. Its core contribution is the explicit treatment of memory
retention as a dynamic process rather than static storage.

Titans introduces a neural long-term memory module that complements
short-term attention, framing attention as short-term memory and
persistent neural memory as a longer-lived store. The later MIRAS
framework generalizes this view by casting sequence models as
associative memory systems with explicit retention and attentional bias
mechanisms.

These approaches motivate the idea that memory should not be identified
with the context window alone. Instead, memory should be structured,
persistent, and updated through principled retention mechanisms.

\subsection{Episodic segmentation and
surprise}

EM-LLM shows that long-context performance benefits from segmenting
experience into events using Bayesian surprise and graph-theoretic
boundary refinement. This is highly relevant to systems that must
distinguish routine background information from important shifts in
world state.

In cognitive science and neuroscience, prediction error has long been
linked to memory updating, event segmentation, and adaptive learning.
Unexpected events can either strengthen encoding or trigger the
formation of a new memory boundary. This motivates the use of surprise
as a write-time salience signal in memory systems.

\subsection{Graph-based memory and
retrieval}

GraphRAG and HippoRAG demonstrate that graph structure can improve
retrieval beyond flat vector similarity. GraphRAG builds entity-centric
knowledge graphs and supports more global, query-focused reasoning over
corpora. HippoRAG combines graph structure with personalized PageRank to
mimic aspects of long-term associative recall.

These systems support the Neocortex design choice to maintain a
coherence graph over entities, relations, mentions, and chunks, rather
than relying only on chunk embeddings.

\subsection{Forgetting and
retention}

The Ebbinghaus forgetting curve remains a foundational model for memory
decay, though later work suggests that the precise functional form of
forgetting can vary. For AI systems, the key principle is operational:
memory strength should decrease with age unless reinforced, while
repeated retrieval or interaction should slow future decay.

\section{Problem Formulation}

We consider a language system that must answer user queries over an
evolving corpus of documents and interactions. Let the memory store at
time \(t\) be

\[
\mathcal{M}_t = \{ \mathcal{C}_t, \mathcal{G}_t, \mathcal{L}_t, \mathcal{I}_t \},
\]

where:

\begin{itemize}
\tightlist
\item
  \(\mathcal{C}_t\) is a chunk store with text and embeddings,
\item
  \(\mathcal{G}_t\) is a coherence graph over entities, relations, and
  mentions,
\item
  \(\mathcal{L}_t\) is an append-only state event ledger,
\item
  \(\mathcal{I}_t\) is an interaction history over users, workspaces,
  and entities.
\end{itemize}

The system must support two broad classes of queries:

\begin{enumerate}
\def\labelenumi{\arabic{enumi}.}
\tightlist
\item
  \textbf{Semantic recall queries}, which benefit from retrieving
  relevant chunks, entities, and relations.
\item
  \textbf{Deterministic state queries}, which require reconstruction of
  ordered facts such as current holder, prior holder, location,
  movement, or transfer history.
\end{enumerate}

The challenge is to retrieve a compact, high-value context under a token
budget while preserving correctness and supporting continual updates.

\section{Neocortex Architecture}

\subsection{End-to-end design}

Neocortex continuously performs two processes:

\begin{enumerate}
\def\labelenumi{\arabic{enumi}.}
\tightlist
\item
  \textbf{Ingestion} of new experience into structured memory.
\item
  \textbf{Recall} of context for downstream response generation.
\end{enumerate}

The architecture separates write-time structuring from recall-time
retrieval. This separation is important because many retrieval failures
originate upstream: if memory is not encoded with entity links, state
events, or salience metadata, it cannot later be selectively recalled.

\subsection{Memory layers}

Neocortex contains four complementary memory layers.

\paragraph{Coherence memory}\label{coherence-memory}

Coherence memory stores semantic and relational context over documents,
chunks, entities, and relations. It supports open-ended retrieval where
the answer may require broad contextual grounding rather than a single
explicit state fact.

\paragraph{State ledger memory}\label{state-ledger-memory}

The state ledger is an append-only event stream derived from chunk text
and relation extraction. It records explicit transitions such as
movement, possession, handoff, and spatial relations. This layer is used
for deterministic state reconstruction.

\paragraph{Interaction memory}\label{interaction-memory}

Interaction memory captures user or workspace engagement with entities
and documents. Interaction strength is encoded through actions such as
\texttt{view}, \texttt{read}, \texttt{react}, \texttt{engage}, and
\texttt{create}. Retrieved items are reinforced through access count and
last-accessed metadata.

\paragraph{Forgetting and surprise
dynamics}\label{forgetting-and-surprise-dynamics}

A retention model reduces the salience of stale, unreinforced memory. A
surprise model boosts non-redundant, behavior-changing content relative
to routine or repeated facts.

\section{Ingestion Pipeline}

\subsection{Parse and chunk}

Incoming documents are first normalized and partitioned into bounded
chunks. Chunking is necessary for embedding, retrieval, citation, and
provenance tracking. Each chunk preserves local order and source
metadata.

\subsection{Entity and relation
extraction}

From each chunk, Neocortex extracts entities and relations. These are
linked into a coherence graph with mention-level metadata such as
frequency, position, and source support. Embeddings are stored for
chunks, entities, and possibly relations.

\subsection{Graph persistence}

The system upserts graph structure and chunk metadata into graph/vector
storage. This enables retrieval paths that use both semantic similarity
and graph connectivity.

\subsection{State event ledger
construction}

In parallel with graph persistence, the system appends state events to a
ledger. Each event contains:

\begin{itemize}
\tightlist
\item
  source document and chunk identifiers,
\item
  sentence order and optional timestamps,
\item
  normalized event type,
\item
  involved entities,
\item
  qualifiers and confidence,
\item
  summary text for inspection.
\end{itemize}

This makes state-style queries auditable and order-sensitive.

\subsection{Surprise scoring}

Optionally, each new chunk is compared against a baseline neighborhood
in memory. If the mismatch is high, the system assigns a larger
\texttt{prediction\_error}; if the new information contradicts prior
memory, the associated reward can be negative. This value is stored as
\texttt{reward\_weight} and later used to modulate salience during
retrieval.

\section{Recall and Routing}

\subsection{Query routing}

Given a query \(q\), Neocortex first decides whether the question is
best handled by:

\begin{itemize}
\tightlist
\item
  a \textbf{semantic retrieval path}, or
\item
  a \textbf{state resolver path}.
\end{itemize}

Queries that ask for current location, previous location, holder,
transfer history, or spatial status are routed to the state resolver.
Queries asking for themes, context, explanations, or broader summaries
are routed to semantic retrieval.

\subsection{Semantic recall}

Semantic recall retrieves chunks, entities, and relations using a
blended score:

\[
\text{Score}(m \mid q) =
\text{Rel}(m, q)
\cdot
\text{Retention}(m)
\cdot
\text{InteractionBias}(m)
\cdot
\text{SurpriseWeight}(m),
\]

where \(\text{Rel}(m, q)\) is semantic relevance,
\(\text{Retention}(m)\) captures recency and reinforcement,
\(\text{InteractionBias}(m)\) captures user or workspace importance, and
\(\text{SurpriseWeight}(m)\) prioritizes informative deltas.

\subsection{Deterministic state
resolution}

For state-style questions, the resolver reconstructs answers from
ordered events in the ledger. This can bypass free-form generation when
a deterministic answer exists, reducing hallucination risk and token
usage.

\section{Forgetting and
Reinforcement}

Neocortex models retention using an Ebbinghaus-style decay process. Let
\(a_m\) denote the age in days of memory item \(m\), and let \(n_m\)
denote its reinforcement count. Stability is defined as

\[
S_m = S_0 \left( 1 + \gamma \ln(1 + n_m) \right),
\]

where \(S_0\) is the base stability and \(\gamma\) is a growth factor.
Retention is then

\[
R_m = \exp(-a_m / S_m).
\]

A configurable floor can clamp very low values to zero.

This formalism captures two desired behaviors:

\begin{enumerate}
\def\labelenumi{\arabic{enumi}.}
\tightlist
\item
  recently accessed memory decays more slowly;
\item
  stale, unreinforced memory contributes less unless reactivated.
\end{enumerate}

Each successful retrieval can update \texttt{last\_accessed\_at} and
increment \texttt{access\_count}, creating a closed loop:

\[
\text{interact} \rightarrow \text{recall} \rightarrow \text{reinforce}.
\]

\section{Surprise-Weighted
Memory}

Pure recency and frequency are insufficient because they overvalue
repeated background facts. Neocortex therefore incorporates a
surprise-weighted factor derived from prediction error.

Let \(m\) be a newly ingested chunk and \(\mathcal{N}(m)\) its nearest
prior memory neighborhood. A prediction-error proxy can be defined as
the mismatch between the new chunk and its expected baseline:

\[
\delta_m = 1 - \text{sim}(m, \mathcal{N}(m)),
\]

optionally augmented with graph-overlap disagreement or contradiction
penalties. The stored salience factor can then be written as

\[
W_m = f(\delta_m, c_m),
\]

where \(c_m\) captures contradiction or reward sign. High-surprise
content receives a larger weight during future retrieval, while
contradictory information can either suppress older beliefs or trigger
state updates in the ledger.

Conceptually, this mechanism promotes memory that changes the
system' s world model rather than memory that merely
repeats what is already known.

\section{Why a State Ledger
Matters}

Many practical questions are poorly served by flat retrieval. Consider
queries such as:

\begin{itemize}
\tightlist
\item
  Who currently holds asset \(x\)?
\item
  Where was object \(y\) before it moved to room \(z\)?
\item
  Which entity interacted most recently with node \(u\)?
\end{itemize}

These are not just retrieval tasks. They require ordered state
reconstruction over transitions. By maintaining an append-only ledger,
Neocortex supports deterministic resolution, provenance inspection, and
conflict handling.

This design also separates two forms of truth:

\begin{itemize}
\tightlist
\item
  \textbf{coherence truth}, which emerges from broad supporting context;
\item
  \textbf{state truth}, which is reconstructed from ordered transitions.
\end{itemize}

That distinction improves interpretability and reduces over-reliance on
free-form synthesis.

\section{Token Efficiency and Product
Implications}

The core product claim of Neocortex is that better memory is also better
token economy. Systems consume too many tokens not only because context
windows are small, but because they lack mechanisms to suppress noise
and preserve only what remains behaviorally useful.

Neocortex improves token efficiency in three ways:

\begin{enumerate}
\def\labelenumi{\arabic{enumi}.}
\tightlist
\item
  \textbf{Structured write-time compression} through chunking, entity
  extraction, and event normalization.
\item
  \textbf{Selective recall} through retention-aware, interaction-aware,
  and surprise-aware retrieval.
\item
  \textbf{Deterministic bypass} for state-style questions that do not
  require broad generative reasoning.
\end{enumerate}

This reframes memory from a storage problem into a \textbf{salience
allocation problem}.

\section{Limitations and Open
Questions}

Neocortex is a systems architecture rather than a single end-to-end
trained model, and several issues remain open.

First, surprise estimation is only as good as the baseline against which
novelty is measured. Poor nearest-neighbor memory can misestimate
importance.

Second, event extraction for the state ledger may introduce schema
errors or miss implicit transitions.

Third, retention parameters such as the base stability and growth factor
are application-dependent and may require calibration.

Fourth, user interaction is an imperfect signal of importance.
Frequently accessed information is not always the most correct or most
valuable.

These challenges suggest future work on learned routing, contradiction
management, uncertainty estimation, and offline memory consolidation.

\section{Conclusion}

We presented Neocortex, a high-level memory pipeline for long-term AI
systems that integrates coherence memory, state-ledger memory,
interaction reinforcement, and forgetting plus surprise dynamics. The
architecture is motivated by both recent memory-augmented language-model
research and cognitive theories of retention and prediction error. The
central idea is that memory should be written in structured form,
recalled through routed mechanisms, strengthened through use, and
allowed to forget when it no longer matters. We argue that such systems
are better aligned with both practical retrieval needs and efficient
token use, because they prioritize what is recent, reinforced, and
informative rather than treating all stored text as equally valuable.

\bibliographystyle{unsrt}
\nocite{
  zhong2023memorybank,
  behrouz2025titans,
  behrouz2025allconnected,
  fountas2024emllm,
  jimenez2024hipporag,
  edge2024graphrag,
  wu2025survey,
  lewandowsky2015ebbinghaus,
  clewett2024predictions,
  kim2024predictionerror
}
\bibliography{references}

\end{document}
